%=============================================================================
%  14_deployment.tex  -  Chapter 13: Configuration and Deployment
%=============================================================================
\chapter{Configuration and Deployment}
\label{ch:deploy}

Universal Analyst is designed to run on a cloud GPU box with the language model
local to the machine, and to be brought up with a single notebook or a single
\code{make} command. This chapter covers configuration, the three deployment
shapes, and the developer tooling.

\section{Configuration}
All settings are read through \code{pydantic-settings} from environment
variables (or a \code{.env} file), with sensible defaults and alias support so
the same setting can be spelled several ways. Table~\ref{tab:env} lists the
important keys; the full template ships as \code{.env.example} and is reproduced
in Appendix~\ref{app:env}.

\begin{table}[h]
\centering
\caption{Principal configuration variables.}
\label{tab:env}
\small
\begin{tabularx}{\linewidth}{@{}L{4.3cm} L{3.0cm} X@{}}
\toprule
\textbf{Variable} & \textbf{Default} & \textbf{Purpose} \\
\midrule
\code{OLLAMA\_HOST} & \code{localhost:11434} & Ollama runtime URL. \\
\code{MODEL\_NAME} & \code{gemma4:12b} & Primary answering model. \\
\code{ROUTER\_MODEL\_NAME} & \code{gemma4:e4b} & Fast routing model. \\
\code{EMBEDDING\_MODEL} & \code{all-MiniLM-L6-v2} & RAG embeddings. \\
\code{MAX\_CONTEXT\_TOKENS} & \code{6000} & RAG context budget. \\
\code{CONTEXT\_WINDOW\_TOKENS} & \code{256000} & Front-end meter window. \\
\code{IMAGE\_CHAT\_ENABLED} & \code{true} & Multimodal image path. \\
\code{MAX\_CHAT\_IMAGES} & \code{3} & Images per turn. \\
\code{MAX\_CHAT\_IMAGE\_MB} & \code{5} & Size per image. \\
\code{CHROMA\_PERSIST\_DIRECTORY} & \code{./chroma\_data} & Vector-store path. \\
\code{ALLOWED\_ORIGINS} & \code{localhost:5173,...} & CORS allowlist. \\
\bottomrule
\end{tabularx}
\end{table}

\begin{lstlisting}[style=py,caption={Settings with multi-alias support (\code{app/config.py}).}]
class Settings(BaseSettings):
    model_name: str = "gemma4:12b"
    router_model_name: str = "gemma4:e4b"
    embedding_model: str = "all-MiniLM-L6-v2"
    max_context_tokens: int = 6000
    environment: str = Field(default="development",
        validation_alias=AliasChoices("ENVIRONMENT", "ENV", "environment", "env"))
    # ... image limits, chroma dir, allowed_origins ...
    model_config = SettingsConfigDict(env_file=".env", extra="ignore",
                                      populate_by_name=True)
\end{lstlisting}

\section{Cloud notebook setup}
The primary path for graders and cloud users is \code{main.ipynb}. Running all
cells installs dependencies idempotently, pulls the required Gemma tags
(\code{gemma4:12b} and \code{gemma4:e4b}), and starts both services. If a
required model tag is unavailable, setup \emph{stops with a real error} rather
than degrading silently --- so the runtime can be fixed directly. The forwarded
port 5173 serves the app and 8000 exposes the API docs.

\section{Docker Compose}
Two compose files cover development and production.

\begin{description}[style=nextline,leftmargin=0pt]
  \item[\textcolor{uaIndigoD}{Development (\code{docker-compose.yml})}]
    Bind-mounts the backend and frontend source for live reload, runs the Vite
    dev server directly from a Node image (so hot-module reload and the
    \code{/api} proxy work out of the box), and reaches an Ollama running on the
    host via \code{host.docker.internal}.
  \item[\textcolor{uaIndigoD}{Production (\code{docker-compose.prod.yml})}]
    Builds optimized images and serves the front end's static assets behind
    nginx, which also proxies \code{/api} to the backend.
\end{description}

\begin{lstlisting}[style=bash,caption={Developer entry points via the Makefile.}]
make up          # dev: backend :8000 + Vite dev server :5173 (Ollama on host)
make prod-up     # prod: built images behind nginx
make down        # stop
make logs        # follow logs
make test        # backend pytest
make lint        # ruff check
make format      # ruff format
\end{lstlisting}

\section{Container images}
The backend \code{Dockerfile} builds a Python\,3.12 service image; the frontend
\code{Dockerfile} builds the production nginx image serving the compiled SPA. The
Makefile detects Docker Compose v2 (\code{docker compose}) and falls back to the
legacy v1 binary if only that is present, so the same commands work across
environments.

\section{Manual startup}
For a bare host, the sequence is: pull the models
(\code{scripts/setup\_ollama.sh}), install backend requirements and launch
\code{uvicorn} on port 8000, then install and run the Vite dev server. Copying
\code{.env.example} to \code{.env} configures models, the Ollama host, the RAG
budget, and CORS.

\begin{lstlisting}[style=bash,caption={Manual startup on a bare GPU host.}]
bash scripts/setup_ollama.sh
cd backend && pip install -r requirements.txt
uvicorn app.main:app --host 0.0.0.0 --port 8000
cd ../frontend && npm install && npm run dev -- --host
\end{lstlisting}

\begin{uawarn}[Ollama placement matters]
Under Docker, the backend reaches Ollama on the \emph{host} through
\code{host.docker.internal}, not \code{localhost}. Getting this wrong is the most
common local-setup mistake; the environment template documents both forms.
\end{uawarn}
